% GENERATED FILE. Edit canonical lessons, then run npm run generate:books.
% canonical-source-hash: fnv1a64:201cd13fe56c456e
% canonical-lessons: SA-C510-avakarah, SA-C510-utsarayati, SA-C510-vyavasthapayati, SA-C510-punjah, SA-C510-sammarjanam

\chapter{Clearing the Clutter}
\label{ch:sa-clearing-the-clutter}
\hlchaptermodality{510}

\begin{chapteropening}
\textbf{By the end of this chapter:} \emph{I can talk about rubbish, sweeping out and putting things in order.}

Complete the last lesson of chapter 510: I can talk about rubbish, sweeping out and putting things in order.
\end{chapteropening}

\section[\texorpdfstring{\sk{अवकरः} (avakaraḥ)}{अवकरः (avakaraḥ)}]{\sk{अवकरः} (avakaraḥ) — rubbish; sweepings}

\label{lesson:SA-C510-avakarah}



\begin{quote}
Before the new one: say the Sanskrit for trickles, then the Sanskrit for a repair.
\end{quote}

\subsection*{You'll want to know: \sk{अवकरः}}
\textbf{\sk{अवकरः}} — \emph{avakaraḥ} — \textquotedblleft{}rubbish; sweepings\textquotedblright{}.

\begin{quote}
\sk{प्रकोष्ठे} \sk{बहु} \sk{अवकरः} \sk{अस्ति।}
\end{quote}

\begin{quote}
\emph{prakoṣṭhe bahu avakaraḥ asti.}
\end{quote}

\begin{quote}
There is a lot of rubbish in the room.
\end{quote}

A masculine noun.

\begin{cousinweb}
\textbf{\sk{अवकरः}} (\emph{avakaraḥ}) is built on \emph{ava}, \textquotedblleft{}down\textquotedblright{}, and the root √kṝ, \textquotedblleft{}to scatter\textquotedblright{}; MW gives \emph{avakara} as dust or sweepings (MW).
\end{cousinweb}

\begin{grammarlens}[title={What you've built}]
One more word for describing what you lived and felt.
\end{grammarlens}

\subsection*{Your turn}
\begin{itemize}
\raggedright
  \item \emph{Say it:} \emph{avakaraḥ}
  \item \emph{Say it:} \emph{avakaraḥ}, once more
  \item \emph{Say it:} \emph{prakoṣṭhe bahu avakaraḥ asti.} — There is a lot of rubbish in the room.
  \item \emph{Recall:} say the Sanskrit for trickles, then the Sanskrit for a repair, then say \emph{avakaraḥ} again
\end{itemize}

\begin{tcolorbox}[breakable,colback=teal!4,colframe=teal!35!black,title={Before you move on}]
What does \sk{अवकरः} mean? (\textquotedblleft{}Rubbish; sweepings\textquotedblright{}.) Say it once more.
\end{tcolorbox}
\section[\texorpdfstring{\sk{उत्सारयति} (utsārayati)}{उत्सारयति (utsārayati)}]{\sk{उत्सारयति} (utsārayati) — clears away}

\label{lesson:SA-C510-utsarayati}



\begin{quote}
Before the new one: say the Sanskrit for a repair, then the Sanskrit for rubbish.
\end{quote}

\subsection*{You'll want to know: \sk{उत्सारयति}}
\textbf{\sk{उत्सारयति}} — \emph{utsārayati} — \textquotedblleft{}clears away\textquotedblright{}.

\begin{quote}
\sk{माता} \sk{अवकरं} \sk{सम्मार्जन्या} \sk{उत्सारयति।}
\end{quote}

\begin{quote}
\emph{mātā avakaraṃ sammārjanyā utsārayati.}
\end{quote}

\begin{quote}
Mother clears away the rubbish with a broom.
\end{quote}

Past (she …): \sk{सा} \sk{अवकरम्} \sk{उत्सारितवती}, \emph{sā avakaram utsāritavatī}.

\begin{cousinweb}
\textbf{\sk{उत्सारयति}} (\emph{utsārayati}) is built on \emph{ud}, \textquotedblleft{}out\textquotedblright{}, and the causative of the root √sṛ, \textquotedblleft{}to run, to flow\textquotedblright{}; MW gives \emph{utsārayati} as to drive away, to put or throw away (MW).
\end{cousinweb}

\begin{grammarlens}[title={What you've built}]
One more word for describing what you lived and felt.
\end{grammarlens}

\subsection*{Your turn}
\begin{itemize}
\raggedright
  \item \emph{Say it:} \emph{utsārayati}
  \item \emph{Say it:} \emph{utsārayati}, once more
  \item \emph{Say it:} \emph{mātā avakaraṃ sammārjanyā utsārayati.} — Mother clears away the rubbish with a broom.
  \item \emph{Recall:} say the Sanskrit for a repair, then the Sanskrit for rubbish, then say \emph{utsārayati} again
\end{itemize}

\begin{tcolorbox}[breakable,colback=teal!4,colframe=teal!35!black,title={Before you move on}]
What does \sk{उत्सारयति} mean? (\textquotedblleft{}Clears away\textquotedblright{}.) Say it once more.
\end{tcolorbox}
\section[\texorpdfstring{\sk{व्यवस्थापयति} (vyavasthāpayati)}{व्यवस्थापयति (vyavasthāpayati)}]{\sk{व्यवस्थापयति} (vyavasthāpayati) — puts in order}

\label{lesson:SA-C510-vyavasthapayati}



\begin{quote}
Before the new one: say the Sanskrit for rubbish, then the Sanskrit for clears away.
\end{quote}

\subsection*{You'll want to know: \sk{व्यवस्थापयति}}
\textbf{\sk{व्यवस्थापयति}} — \emph{vyavasthāpayati} — \textquotedblleft{}puts in order\textquotedblright{}.

\begin{quote}
\sk{सः} \sk{पुस्तकानि} \sk{व्यवस्थापयति।}
\end{quote}

\begin{quote}
\emph{saḥ pustakāni vyavasthāpayati.}
\end{quote}

\begin{quote}
He puts the books in order.
\end{quote}

Past (he …): \sk{सः} \sk{पुस्तकानि} \sk{व्यवस्थापितवान्}, \emph{saḥ pustakāni vyavasthāpitavān}.

\begin{cousinweb}
\textbf{\sk{व्यवस्थापयति}} (\emph{vyavasthāpayati}) is built on \emph{vi}, \textquotedblleft{}apart\textquotedblright{}, \emph{ava}, \textquotedblleft{}down\textquotedblright{}, and the causative of the root √sthā, \textquotedblleft{}to stand\textquotedblright{}: to set each thing down in its place; MW gives the causative \emph{vyavasthāpayati} as to put down, to place, to fix (MW).
\end{cousinweb}

\begin{grammarlens}[title={What you've built}]
One more word for describing what you lived and felt.
\end{grammarlens}

\subsection*{Your turn}
\begin{itemize}
\raggedright
  \item \emph{Say it:} \emph{vyavasthāpayati}
  \item \emph{Say it:} \emph{vyavasthāpayati}, once more
  \item \emph{Say it:} \emph{saḥ pustakāni vyavasthāpayati.} — He puts the books in order.
  \item \emph{Recall:} say the Sanskrit for rubbish, then the Sanskrit for clears away, then say \emph{vyavasthāpayati} again
\end{itemize}

\begin{tcolorbox}[breakable,colback=teal!4,colframe=teal!35!black,title={Before you move on}]
What does \sk{व्यवस्थापयति} mean? (\textquotedblleft{}Puts in order\textquotedblright{}.) Say it once more.
\end{tcolorbox}
\section[\texorpdfstring{\sk{पुञ्जः} (puñjaḥ)}{पुञ्जः (puñjaḥ)}]{\sk{पुञ्जः} (puñjaḥ) — a pile; a heap}

\label{lesson:SA-C510-punjah}



\begin{quote}
Before the new one: say the Sanskrit for clears away, then the Sanskrit for puts in order.
\end{quote}

\subsection*{You'll want to know: \sk{पुञ्जः}}
\textbf{\sk{पुञ्जः}} — \emph{puñjaḥ} — \textquotedblleft{}a pile; a heap\textquotedblright{}.

\begin{quote}
\sk{गृहस्य} \sk{बहिः} \sk{अवकरस्य} \sk{पुञ्जः} \sk{अस्ति।}
\end{quote}

\begin{quote}
\emph{gṛhasya bahiḥ avakarasya puñjaḥ asti.}
\end{quote}

\begin{quote}
There is a heap of rubbish outside the house.
\end{quote}

A masculine noun.

\begin{cousinweb}
\textbf{\sk{पुञ्जः}} (\emph{puñjaḥ}) is the word MW records for a heap, a mass, a multitude; MW gives no derivation, so its origin is uncertain (MW).
\end{cousinweb}

\begin{grammarlens}[title={What you've built}]
One more word for describing what you lived and felt.
\end{grammarlens}

\subsection*{Your turn}
\begin{itemize}
\raggedright
  \item \emph{Say it:} \emph{puñjaḥ}
  \item \emph{Say it:} \emph{puñjaḥ}, once more
  \item \emph{Say it:} \emph{gṛhasya bahiḥ avakarasya puñjaḥ asti.} — There is a heap of rubbish outside the house.
  \item \emph{Recall:} say the Sanskrit for clears away, then the Sanskrit for puts in order, then say \emph{puñjaḥ} again
\end{itemize}

\begin{tcolorbox}[breakable,colback=teal!4,colframe=teal!35!black,title={Before you move on}]
What does \sk{पुञ्जः} mean? (\textquotedblleft{}A pile; a heap\textquotedblright{}.) Say it once more.
\end{tcolorbox}
\section[\texorpdfstring{\sk{सम्मार्जनम्} (sammārjanam)}{सम्मार्जनम् (sammārjanam)}]{\sk{सम्मार्जनम्} (sammārjanam) — a thorough cleaning}

\label{lesson:SA-C510-sammarjanam}



\begin{quote}
Before the new one: say the Sanskrit for puts in order, then the Sanskrit for a pile.
\end{quote}

\subsection*{You'll want to know: \sk{सम्मार्जनम्}}
\textbf{\sk{सम्मार्जनम्}} — \emph{sammārjanam} — \textquotedblleft{}a thorough cleaning\textquotedblright{}.

\begin{quote}
\sk{उत्सवात्} \sk{पूर्वं} \sk{गृहस्य} \sk{सम्मार्जनं} \sk{भवति।}
\end{quote}

\begin{quote}
\emph{utsavāt pūrvaṃ gṛhasya sammārjanaṃ bhavati.}
\end{quote}

\begin{quote}
Before the festival the house gets a thorough cleaning.
\end{quote}

A neuter noun.

\begin{cousinweb}
\textbf{\sk{सम्मार्जनम्}} (\emph{sammārjanam}) is built on \emph{sam}, \textquotedblleft{}thoroughly\textquotedblright{}, and the root √mṛj, \textquotedblleft{}to wipe, to clean\textquotedblright{}, which also gives \emph{mārjayati}; MW gives \emph{sammārjana} as sweeping or cleansing thoroughly (MW).
\end{cousinweb}

\begin{grammarlens}[title={What you've built}]
One more word for describing what you lived and felt.
\end{grammarlens}

\subsection*{Your turn}
\begin{itemize}
\raggedright
  \item \emph{Say it:} \emph{sammārjanam}
  \item \emph{Say it:} \emph{sammārjanam}, once more
  \item \emph{Say it:} \emph{utsavāt pūrvaṃ gṛhasya sammārjanaṃ bhavati.} — Before the festival the house gets a thorough cleaning.
  \item \emph{Recall:} say the Sanskrit for puts in order, then the Sanskrit for a pile, then say \emph{sammārjanam} again
\end{itemize}

\begin{tcolorbox}[breakable,colback=teal!4,colframe=teal!35!black,title={Before you move on}]
What does \sk{सम्मार्जनम्} mean? (\textquotedblleft{}A thorough cleaning\textquotedblright{}.) Say it once more.
\end{tcolorbox}
